% 2026-10-08 parked (lean set; replaced by an MCS over the features): the whole subsection; its \input in appendix_campaign_16h.tex is parked too
% % ============================================================================
% % 2026-10-06 (user decision: "make the lean set the third input set in the paper, but it has
% % to be a follow up to the feature importance study").  A new appendix subsection after the
% % 16:00 campaign (Appendix D.19), which keeps its own framing: what the importance study of
% % D.19 finds, the lean set defined from it, its evaluation against HAR + calendar and all
% % features, and the selection caveat.  Every number is a \fl... macro of
% % generated/featimp_lean_numbers.tex (writeup/make_featimp_lean_tex.py); the two tables are
% % generated/featimp_lean_{importance,results}.tex.  Display names carry no adj_ prefix.
% % Input from sections/appendix_campaign_16h.tex, at its end.
% % ============================================================================
% \ifdefined\flPreRefits\else% AUTO-GENERATED by writeup/make_featimp_lean_tex.py -- do not edit.
% The lean set (lean_vol_vix) as a follow-up to the feature-importance study (2026-10-06):
% importance ranks (results/feature_importance_1530_dedup/, all_features, series level) and
% lean vs all_features vs baseline over the paper model set (point estimates; DM / HAC t).
\newcommand{\flPreRefits}{39}
\newcommand{\flPreTailDays}{390}
\newcommand{\flPreRefitsEighteen}{13}
\newcommand{\flPreTailDaysEighteen}{130}
\newcommand{\flPreRefitsNineteen}{26}
\newcommand{\flPreTailDaysNineteen}{260}
\newcommand{\flPreTailFirst}{2018-06-25}
\newcommand{\flPreTailLast}{2020-01-14}
\newcommand{\flPreTiedCells}{8}
\newcommand{\flCells}{37}
\newcommand{\flSeriesN}{49}
\newcommand{\flModelsN}{5}
\newcommand{\flTop}{10}
\newcommand{\flRankMedHar}{1}
\newcommand{\flTopTenHar}{35}
\newcommand{\flRankBestHar}{1}
\newcommand{\flPlaceHar}{1}
\newcommand{\flPreRankHar}{1}
\newcommand{\flPreTopTenHar}{37}
\newcommand{\flRankMedSumret}{3}
\newcommand{\flTopTenSumret}{32}
\newcommand{\flRankBestSumret}{2}
\newcommand{\flPlaceSumret}{2}
\newcommand{\flPreRankSumret}{3}
\newcommand{\flPreTopTenSumret}{34}
\newcommand{\flRankMedSumabsret}{5}
\newcommand{\flTopTenSumabsret}{28}
\newcommand{\flRankBestSumabsret}{1}
\newcommand{\flPlaceSumabsret}{3}
\newcommand{\flPreRankSumabsret}{4}
\newcommand{\flPreTopTenSumabsret}{30}
\newcommand{\flRankMedSumretthree}{20}
\newcommand{\flTopTenSumretthree}{8}
\newcommand{\flRankBestSumretthree}{3}
\newcommand{\flPlaceSumretthree}{14}
\newcommand{\flPreRankSumretthree}{18}
\newcommand{\flPreTopTenSumretthree}{8}
\newcommand{\flRankMedSumretfour}{18}
\newcommand{\flTopTenSumretfour}{11}
\newcommand{\flRankBestSumretfour}{3}
\newcommand{\flPlaceSumretfour}{10}
\newcommand{\flPreRankSumretfour}{13}
\newcommand{\flPreTopTenSumretfour}{16}
\newcommand{\flRankMedSumprettwo}{33}
\newcommand{\flTopTenSumprettwo}{2}
\newcommand{\flRankBestSumprettwo}{9}
\newcommand{\flPlaceSumprettwo}{40}
\newcommand{\flPreRankSumprettwo}{26}
\newcommand{\flPreTopTenSumprettwo}{0}
\newcommand{\flRankMedSumbipow}{33}
\newcommand{\flTopTenSumbipow}{1}
\newcommand{\flRankBestSumbipow}{9}
\newcommand{\flPlaceSumbipow}{39}
\newcommand{\flPreRankSumbipow}{25}
\newcommand{\flPreTopTenSumbipow}{4}
\newcommand{\flRankMedSumautocov}{12}
\newcommand{\flTopTenSumautocov}{15}
\newcommand{\flRankBestSumautocov}{1}
\newcommand{\flPlaceSumautocov}{8}
\newcommand{\flPreRankSumautocov}{12}
\newcommand{\flPreTopTenSumautocov}{17}
\newcommand{\flRankMedSumvolume}{5}
\newcommand{\flTopTenSumvolume}{23}
\newcommand{\flRankBestSumvolume}{3}
\newcommand{\flPlaceSumvolume}{4}
\newcommand{\flPreRankSumvolume}{11}
\newcommand{\flPreTopTenSumvolume}{17}
\newcommand{\flRankMedVix}{19}
\newcommand{\flTopTenVix}{9}
\newcommand{\flRankBestVix}{3}
\newcommand{\flPlaceVix}{12}
\newcommand{\flPreRankVix}{15}
\newcommand{\flPreTopTenVix}{8}
\newcommand{\flOtherOneLabel}{\texttt{calendar (DOW\_*, is\_*, ...)}}
\newcommand{\flRankMedOtherOne}{6}
\newcommand{\flTopTenOtherOne}{24}
\newcommand{\flRankBestOtherOne}{2}
\newcommand{\flPlaceOtherOne}{5}
\newcommand{\flOtherTwoLabel}{\texttt{sumabsret\_vwstock\_ma\_*}}
\newcommand{\flRankMedOtherTwo}{7}
\newcommand{\flTopTenOtherTwo}{27}
\newcommand{\flRankBestOtherTwo}{2}
\newcommand{\flPlaceOtherTwo}{6}
\newcommand{\flOtherThreeLabel}{\texttt{sumabsret\_ewstock\_ma\_*}}
\newcommand{\flRankMedOtherThree}{8}
\newcommand{\flTopTenOtherThree}{21}
\newcommand{\flRankBestOtherThree}{3}
\newcommand{\flPlaceOtherThree}{7}
\newcommand{\flOtherFourLabel}{\texttt{ofi\_ewstock\_ma\_*}}
\newcommand{\flRankMedOtherFour}{12}
\newcommand{\flTopTenOtherFour}{15}
\newcommand{\flRankBestOtherFour}{7}
\newcommand{\flPlaceOtherFour}{9}
\newcommand{\flOtherFiveLabel}{\texttt{ofi\_vwstock\_ma\_*}}
\newcommand{\flRankMedOtherFive}{18}
\newcommand{\flTopTenOtherFive}{5}
\newcommand{\flRankBestOtherFive}{7}
\newcommand{\flPlaceOtherFive}{11}
\newcommand{\flRankMedVixthreem}{21}
\newcommand{\flTopTenVixthreem}{8}
\newcommand{\flRankBestVixthreem}{4}
\newcommand{\flPlaceVixthreem}{17}
\newcommand{\flRankMedStocktwits}{36}
\newcommand{\flTopTenStocktwits}{12}
\newcommand{\flRankBestStocktwits}{2}
\newcommand{\flPlaceStocktwits}{46}
\newcommand{\flLeanInTopTenN}{5}
\newcommand{\flLeanRankMedMax}{19}
\newcommand{\flLeanN}{6}
\newcommand{\flClusterThreeOneTwoFiveCorr}{0.85}
\newcommand{\flClusterOneFiveCorr}{0.82}
\newcommand{\flClusterSixTwentyFiveCorr}{0.89}
\newcommand{\flClusterOneTwentyFiveCorr}{0.89}
\newcommand{\flClusterTwentyFiveCorr}{0.88}
\newcommand{\flClusterN}{5}
\newcommand{\flClusterCorrMin}{0.82}
\newcommand{\flClusterCorrMax}{0.89}
\newcommand{\flClusterLagsN}{6}
\newcommand{\flPreClusterRows}{2382}
\newcommand{\flPreClusterFirst}{2010-07-12}
\newcommand{\flPreClusterLast}{2020-01-02}
\newcommand{\flPreClusterMinCorrLagOne}{0.87}
\newcommand{\flPreClusterMinCorrLagFive}{0.87}
\newcommand{\flPreClusterMinCorrLagTwentyFive}{0.85}
\newcommand{\flPreClusterMinCorrLagOneTwentyFive}{0.84}
\newcommand{\flPreClusterMinCorrLagSixTwentyFive}{0.79}
\newcommand{\flPreClusterMinCorrLagThreeOneTwoFive}{0.75}
\newcommand{\flLeanForecastsN}{51}
\newcommand{\flLinearN}{3}
\newcommand{\flLinearQLean}{0.0937}
\newcommand{\flLinearShLean}{1.88}
\newcommand{\flLinearShXLean}{1.40}
\newcommand{\flLinearQAllF}{0.0942}
\newcommand{\flLinearShAllF}{1.88}
\newcommand{\flLinearShXAllF}{1.40}
\newcommand{\flLinearQBase}{0.0946}
\newcommand{\flLinearShBase}{1.37}
\newcommand{\flLinearShXBase}{0.89}
\newcommand{\flLinearVsAllFQPctMed}{\ensuremath{-}0.5}
\newcommand{\flLinearVsAllFShMed}{+0.01}
\newcommand{\flLinearVsAllFQBetter}{2}
\newcommand{\flLinearVsAllFQBetterSig}{0}
\newcommand{\flLinearVsAllFQWorse}{1}
\newcommand{\flLinearVsAllFQWorseSig}{0}
\newcommand{\flLinearVsAllFShAbove}{2}
\newcommand{\flLinearVsAllFShAboveSig}{0}
\newcommand{\flLinearVsAllFShBelow}{1}
\newcommand{\flLinearVsAllFShBelowSig}{0}
\newcommand{\flLinearVsBaseQPctMed}{\ensuremath{-}1.0}
\newcommand{\flLinearVsBaseShMed}{+0.51}
\newcommand{\flLinearVsBaseQBetter}{2}
\newcommand{\flLinearVsBaseQBetterSig}{0}
\newcommand{\flLinearVsBaseQWorse}{1}
\newcommand{\flLinearVsBaseQWorseSig}{0}
\newcommand{\flLinearVsBaseShAbove}{3}
\newcommand{\flLinearVsBaseShAboveSig}{0}
\newcommand{\flLinearVsBaseShBelow}{0}
\newcommand{\flLinearVsBaseShBelowSig}{0}
\newcommand{\flPooledN}{3}
\newcommand{\flPooledQLean}{0.1038}
\newcommand{\flPooledShLean}{1.23}
\newcommand{\flPooledShXLean}{0.74}
\newcommand{\flPooledQAllF}{0.1021}
\newcommand{\flPooledShAllF}{1.44}
\newcommand{\flPooledShXAllF}{0.96}
\newcommand{\flPooledQBase}{0.1065}
\newcommand{\flPooledShBase}{0.89}
\newcommand{\flPooledShXBase}{0.41}
\newcommand{\flPooledVsAllFQPctMed}{+1.8}
\newcommand{\flPooledVsAllFShMed}{\ensuremath{-}0.21}
\newcommand{\flPooledVsAllFQBetter}{1}
\newcommand{\flPooledVsAllFQBetterSig}{0}
\newcommand{\flPooledVsAllFQWorse}{2}
\newcommand{\flPooledVsAllFQWorseSig}{0}
\newcommand{\flPooledVsAllFShAbove}{1}
\newcommand{\flPooledVsAllFShAboveSig}{0}
\newcommand{\flPooledVsAllFShBelow}{2}
\newcommand{\flPooledVsAllFShBelowSig}{0}
\newcommand{\flPooledVsBaseQPctMed}{\ensuremath{-}2.5}
\newcommand{\flPooledVsBaseShMed}{+0.40}
\newcommand{\flPooledVsBaseQBetter}{3}
\newcommand{\flPooledVsBaseQBetterSig}{0}
\newcommand{\flPooledVsBaseQWorse}{0}
\newcommand{\flPooledVsBaseQWorseSig}{0}
\newcommand{\flPooledVsBaseShAbove}{3}
\newcommand{\flPooledVsBaseShAboveSig}{0}
\newcommand{\flPooledVsBaseShBelow}{0}
\newcommand{\flPooledVsBaseShBelowSig}{0}
\newcommand{\flTreeShipN}{6}
\newcommand{\flTreeShipQLean}{0.0924}
\newcommand{\flTreeShipShLean}{1.47}
\newcommand{\flTreeShipShXLean}{0.98}
\newcommand{\flTreeShipQAllF}{0.0958}
\newcommand{\flTreeShipShAllF}{1.50}
\newcommand{\flTreeShipShXAllF}{1.02}
\newcommand{\flTreeShipQBase}{0.0986}
\newcommand{\flTreeShipShBase}{0.93}
\newcommand{\flTreeShipShXBase}{0.45}
\newcommand{\flTreeShipVsAllFQPctMed}{\ensuremath{-}3.3}
\newcommand{\flTreeShipVsAllFShMed}{+0.02}
\newcommand{\flTreeShipVsAllFQBetter}{6}
\newcommand{\flTreeShipVsAllFQBetterSig}{3}
\newcommand{\flTreeShipVsAllFQWorse}{0}
\newcommand{\flTreeShipVsAllFQWorseSig}{0}
\newcommand{\flTreeShipVsAllFShAbove}{3}
\newcommand{\flTreeShipVsAllFShAboveSig}{1}
\newcommand{\flTreeShipVsAllFShBelow}{3}
\newcommand{\flTreeShipVsAllFShBelowSig}{0}
\newcommand{\flTreeShipVsBaseQPctMed}{\ensuremath{-}6.1}
\newcommand{\flTreeShipVsBaseShMed}{+0.51}
\newcommand{\flTreeShipVsBaseQBetter}{6}
\newcommand{\flTreeShipVsBaseQBetterSig}{5}
\newcommand{\flTreeShipVsBaseQWorse}{0}
\newcommand{\flTreeShipVsBaseQWorseSig}{0}
\newcommand{\flTreeShipVsBaseShAbove}{6}
\newcommand{\flTreeShipVsBaseShAboveSig}{0}
\newcommand{\flTreeShipVsBaseShBelow}{0}
\newcommand{\flTreeShipVsBaseShBelowSig}{0}
\newcommand{\flTreeRsN}{12}
\newcommand{\flTreeRsQLean}{0.0999}
\newcommand{\flTreeRsShLean}{1.71}
\newcommand{\flTreeRsShXLean}{1.23}
\newcommand{\flTreeRsQAllF}{0.0995}
\newcommand{\flTreeRsShAllF}{1.41}
\newcommand{\flTreeRsShXAllF}{0.93}
\newcommand{\flTreeRsQBase}{0.1021}
\newcommand{\flTreeRsShBase}{0.96}
\newcommand{\flTreeRsShXBase}{0.48}
\newcommand{\flTreeRsVsAllFQPctMed}{\ensuremath{-}0.4}
\newcommand{\flTreeRsVsAllFShMed}{+0.18}
\newcommand{\flTreeRsVsAllFQBetter}{6}
\newcommand{\flTreeRsVsAllFQBetterSig}{3}
\newcommand{\flTreeRsVsAllFQWorse}{6}
\newcommand{\flTreeRsVsAllFQWorseSig}{0}
\newcommand{\flTreeRsVsAllFShAbove}{9}
\newcommand{\flTreeRsVsAllFShAboveSig}{0}
\newcommand{\flTreeRsVsAllFShBelow}{3}
\newcommand{\flTreeRsVsAllFShBelowSig}{0}
\newcommand{\flTreeRsVsBaseQPctMed}{\ensuremath{-}2.6}
\newcommand{\flTreeRsVsBaseShMed}{+0.79}
\newcommand{\flTreeRsVsBaseQBetter}{11}
\newcommand{\flTreeRsVsBaseQBetterSig}{0}
\newcommand{\flTreeRsVsBaseQWorse}{1}
\newcommand{\flTreeRsVsBaseQWorseSig}{0}
\newcommand{\flTreeRsVsBaseShAbove}{11}
\newcommand{\flTreeRsVsBaseShAboveSig}{3}
\newcommand{\flTreeRsVsBaseShBelow}{1}
\newcommand{\flTreeRsVsBaseShBelowSig}{0}
\newcommand{\flTreeOptN}{27}
\newcommand{\flTreeOptQLean}{0.0974}
\newcommand{\flTreeOptShLean}{1.30}
\newcommand{\flTreeOptShXLean}{0.81}
\newcommand{\flTreeOptQAllF}{0.1002}
\newcommand{\flTreeOptShAllF}{1.23}
\newcommand{\flTreeOptShXAllF}{0.75}
\newcommand{\flTreeOptQBase}{0.1010}
\newcommand{\flTreeOptShBase}{0.71}
\newcommand{\flTreeOptShXBase}{0.23}
\newcommand{\flTreeOptVsAllFQPctMed}{\ensuremath{-}3.2}
\newcommand{\flTreeOptVsAllFShMed}{+0.18}
\newcommand{\flTreeOptVsAllFQBetter}{22}
\newcommand{\flTreeOptVsAllFQBetterSig}{7}
\newcommand{\flTreeOptVsAllFQWorse}{5}
\newcommand{\flTreeOptVsAllFQWorseSig}{1}
\newcommand{\flTreeOptVsAllFShAbove}{19}
\newcommand{\flTreeOptVsAllFShAboveSig}{1}
\newcommand{\flTreeOptVsAllFShBelow}{8}
\newcommand{\flTreeOptVsAllFShBelowSig}{0}
\newcommand{\flTreeOptVsBaseQPctMed}{\ensuremath{-}3.9}
\newcommand{\flTreeOptVsBaseShMed}{+0.73}
\newcommand{\flTreeOptVsBaseQBetter}{26}
\newcommand{\flTreeOptVsBaseQBetterSig}{3}
\newcommand{\flTreeOptVsBaseQWorse}{1}
\newcommand{\flTreeOptVsBaseQWorseSig}{0}
\newcommand{\flTreeOptVsBaseShAbove}{27}
\newcommand{\flTreeOptVsBaseShAboveSig}{6}
\newcommand{\flTreeOptVsBaseShBelow}{0}
\newcommand{\flTreeOptVsBaseShBelowSig}{0}
\newcommand{\flSetN}{51}
\newcommand{\flSetQLean}{0.0974}
\newcommand{\flSetShLean}{1.40}
\newcommand{\flSetShXLean}{0.90}
\newcommand{\flSetQAllF}{0.0994}
\newcommand{\flSetShAllF}{1.37}
\newcommand{\flSetShXAllF}{0.88}
\newcommand{\flSetQBase}{0.1010}
\newcommand{\flSetShBase}{0.79}
\newcommand{\flSetShXBase}{0.31}
\newcommand{\flSetVsAllFQPctMed}{\ensuremath{-}2.4}
\newcommand{\flSetVsAllFShMed}{+0.12}
\newcommand{\flSetVsAllFQBetter}{37}
\newcommand{\flSetVsAllFQBetterSig}{13}
\newcommand{\flSetVsAllFQWorse}{14}
\newcommand{\flSetVsAllFQWorseSig}{1}
\newcommand{\flSetVsAllFShAbove}{34}
\newcommand{\flSetVsAllFShAboveSig}{2}
\newcommand{\flSetVsAllFShBelow}{17}
\newcommand{\flSetVsAllFShBelowSig}{0}
\newcommand{\flSetVsBaseQPctMed}{\ensuremath{-}3.6}
\newcommand{\flSetVsBaseShMed}{+0.69}
\newcommand{\flSetVsBaseQBetter}{48}
\newcommand{\flSetVsBaseQBetterSig}{8}
\newcommand{\flSetVsBaseQWorse}{3}
\newcommand{\flSetVsBaseQWorseSig}{0}
\newcommand{\flSetVsBaseShAbove}{50}
\newcommand{\flSetVsBaseShAboveSig}{9}
\newcommand{\flSetVsBaseShBelow}{1}
\newcommand{\flSetVsBaseShBelowSig}{0}
\fi
% \subsection{A Lean Input Set, Chosen from the Importance Study}
% \label{sec:app_lean}
% 
% \paragraph{What the importance study finds.} The importance study of
% Appendix~\ref{sec:app_camp_featimp} ranks the \flSeriesN{} series of the
% \texttt{all\_features} design for the per-bar ridge, the lasso and the three trees by
% impurity and split counts, permutation, TreeSHAP and, for the linear models, standardized
% weights and drop-column losses: \flCells{} model-by-measure cells, rank 1 the most important
% series in a cell (Table~\ref{tab:app_lean_importance}). The target's own HAR lags lead: median
% rank \flRankMedHar, in the top \flTop{} of \flTopTenHar{} of the \flCells{} cells. Next come
% the signed return (\texttt{sumret\_ma\_*}, median rank \flRankMedSumret), the absolute return
% (\texttt{sumabsret\_ma\_*}, \flRankMedSumabsret), E-mini volume (\texttt{sumvolume\_ma\_*},
% \flRankMedSumvolume) and the return autocovariance (\texttt{sumautocov\_ma\_*},
% \flRankMedSumautocov). Bipower variation (\texttt{sumbipow\_ma\_*}, median rank
% \flRankMedSumbipow) and the upside squared return (\texttt{sumpret2\_ma\_*},
% \flRankMedSumprettwo) sit in one correlation cluster with the absolute and the quartic return
% (\texttt{sumret4\_ma\_*}, \flRankMedSumretfour) at each of the \flClusterLagsN{} lags, with a
% smallest within-cluster correlation of \flClusterCorrMin{} to \flClusterCorrMax{} across the
% lags, and add nothing beyond them. The returns of the index constituents rank (the absolute
% constituent returns at median ranks \flRankMedOtherTwo{} and \flRankMedOtherThree), but they
% duplicate the E-mini moments.
% % source (no macro): results/feature_importance_1530_dedup/SUMMARY.md -- the VIX level is in
% % the top three of the lasso's permutation and standardized-weight measures and outside the
% % trees' top ranks.
% The VIX level (\texttt{vix\_ma\_*}, median rank \flRankMedVix) matters more for the linear
% models than for the trees.
% 
% \begin{table}[H]
% \centering
% \caption{Importance of the series of the \texttt{all\_features} design over the
% \flCells{} model-by-measure cells of the importance study (rank 1 = most important of
% \flSeriesN{} series): median rank and the number of cells with the series in the top
% \flTop, on all held-out tails (2018--2024) and on the tails of the refits trained before the
% trade period, the best rank, and the place by median rank. Lean: the series of the lean set.}
% \label{tab:app_lean_importance}
% \campfit{\scriptsize% AUTO-GENERATED by writeup/make_featimp_lean_tex.py -- do not edit.
% Source: results/feature_importance_1530_dedup/by_measure/imp_all_features_<model>_<measure>.csv (level 'series', 37 cells: rank 1 = most important among 49 series); pre-2020 = the y2018 / y2019 columns weighted by their tail days, re-ranked
\begin{tabular}{llcrrrrrr}
\toprule
 & & & \multicolumn{2}{c}{all tails, 2018--2024} & \multicolumn{2}{c}{pre-2020 tails} & & \\
\cmidrule(lr){4-5}\cmidrule(lr){6-7}
series & what & lean & \shortstack{median\\rank} & \shortstack{top-10\\cells} & \shortstack{median\\rank} & \shortstack{top-10\\cells} & \shortstack{best\\rank} & \shortstack{place by\\median} \\
\midrule
\texttt{har\_ma\_*} & realized variance (the target's HAR ladder) &  & 1 & 35 & 1 & 37 & 1 & 1 \\
\midrule
\texttt{sumret\_ma\_*} & signed return & \checkmark & 3 & 32 & 3 & 34 & 2 & 2 \\
\texttt{sumabsret\_ma\_*} & absolute return & \checkmark & 5 & 28 & 4 & 30 & 1 & 3 \\
\texttt{sumret3\_ma\_*} & cubed return &  & 20 & 8 & 18 & 8 & 3 & 14 \\
\texttt{sumret4\_ma\_*} & quartic return & \checkmark & 18 & 11 & 13 & 16 & 3 & 10 \\
\texttt{sumpret2\_ma\_*} & upside squared return &  & 33 & 2 & 26 & 0 & 9 & 40 \\
\texttt{sumbipow\_ma\_*} & bipower variation &  & 33 & 1 & 25 & 4 & 9 & 39 \\
\texttt{sumautocov\_ma\_*} & return autocovariance & \checkmark & 12 & 15 & 12 & 17 & 1 & 8 \\
\texttt{sumvolume\_ma\_*} & ES volume & \checkmark & 5 & 23 & 11 & 17 & 3 & 4 \\
\texttt{vix\_ma\_*} & VIX & \checkmark & 19 & 9 & 15 & 8 & 3 & 12 \\
\midrule
\texttt{calendar (DOW\_*, is\_*, ...)} & best-ranked other series &  & 6 & 24 & 6 & 26 & 2 & 5 \\
\texttt{sumabsret\_vwstock\_ma\_*} & best-ranked other series &  & 7 & 27 & 6 & 26 & 2 & 6 \\
\texttt{sumabsret\_ewstock\_ma\_*} & best-ranked other series &  & 8 & 21 & 7 & 21 & 3 & 7 \\
\texttt{ofi\_ewstock\_ma\_*} & best-ranked other series &  & 12 & 15 & 18 & 6 & 7 & 9 \\
\texttt{ofi\_vwstock\_ma\_*} & best-ranked other series &  & 18 & 5 & 20 & 2 & 7 & 11 \\
\texttt{vix3m\_ma\_*} & comparison &  & 21 & 8 & 23 & 2 & 4 & 17 \\
\texttt{stocktwits\_sentiment\_ma\_*} & comparison &  & 36 & 12 & 14 & 11 & 2 & 46 \\
\bottomrule
\end{tabular}
}
% \end{table}
% 
% \paragraph{The lean set.} The lean set, named \texttt{lean\_vol\_vix} in the tables, keeps the
% HAR + calendar backbone of \texttt{baseline} and adds the \flLeanN{} exogenous series that
% reading points to: the signed return (\texttt{sumret}), the absolute return
% (\texttt{sumabsret}), the quartic return (\texttt{sumret4}), the return autocovariance
% (\texttt{sumautocov}), E-mini volume (\texttt{sumvolume}) and the VIX level (\texttt{vix}),
% each through the same ladder of trailing means as every other series. \flLeanInTopTenN{} of
% the \flLeanN{} are among the ten series with the best median rank.
% 
% \paragraph{Its evaluation.} The per-bar linear forecasts, their pooled twins and the trees are
% fitted on the lean set exactly as on \texttt{baseline} and \texttt{all\_features} --- the same
% rows, window,
% refits and tuning --- and scored on the same days, every forecast used as its plain
% back-transform (Table~\ref{tab:app_lean_results}). Over the \flSetN{} forecasts paired across
% the three sets, the median QLIKE is \flSetQLean{} on the lean set, \flSetQAllF{} on
% \texttt{all\_features} and \flSetQBase{} on \texttt{baseline}, and the median
% $\mathrm{sign}(s)$ Sharpe ratio \flSetShLean, \flSetShAllF{} and \flSetShBase. Against
% \texttt{all\_features}, the lean forecast has the lower QLIKE in \flSetVsAllFQBetter{} pairs
% (\flSetVsAllFQBetterSig{} significantly) and the higher in \flSetVsAllFQWorse{}
% (\flSetVsAllFQWorseSig{}), and trades higher in \flSetVsAllFShAbove{} (\flSetVsAllFShAboveSig{}
% significantly) and lower in \flSetVsAllFShBelow{} (\flSetVsAllFShBelowSig); the median
% differences are \flSetVsAllFQPctMed\% in QLIKE and \flSetVsAllFShMed{} in the Sharpe ratio.
% Against \texttt{baseline} it has the lower QLIKE in \flSetVsBaseQBetter{}
% (\flSetVsBaseQBetterSig{} significantly) and trades higher in \flSetVsBaseShAbove{}
% (\flSetVsBaseShAboveSig{} significantly), at median differences of \flSetVsBaseQPctMed\% and
% \flSetVsBaseShMed. The pairs share the data and are not independent evidence.
% 
% \begin{table}[H]
% \centering
% \caption{The lean set against \texttt{all\_features} and \texttt{baseline}, per family of the
% model set, every forecast used as its plain back-transform: median QLIKE and median
% $\mathrm{sign}(s)$ Sharpe ratio (midpoint) on each set, and the number of pairs in which the
% lean forecast has the lower QLIKE and the higher Sharpe ratio, with the number significant at
% 5\,\% in parentheses (Diebold--Mariano for QLIKE, the Newey--West $t$ of the daily
% midpoint-return difference for the Sharpe ratio). Point estimates.}
% \label{tab:app_lean_results}
% \campfit{\scriptsize% AUTO-GENERATED by writeup/make_featimp_lean_tex.py -- do not edit.
% Source: results/close_master_table/master_table.csv
% Source: results/close_master_table/master_table_daily.parquet (DM of the daily QLIKE difference, HAC t of the daily mid-return difference)
\begin{tabular}{lrrrrrrrrrrr}
\toprule
 & & \multicolumn{3}{c}{QLIKE (median)} & \multicolumn{3}{c}{Sharpe mid (median)} & \multicolumn{2}{c}{lean $-$ \texttt{all\_features}} & \multicolumn{2}{c}{lean $-$ \texttt{baseline}} \\
\cmidrule(lr){3-5}\cmidrule(lr){6-8}\cmidrule(lr){9-10}\cmidrule(lr){11-12}
family & pairs & lean & all & base & lean & all & base & \shortstack{QLIKE lower\\(DM sig.)} & \shortstack{Sharpe higher\\(HAC sig.)} & \shortstack{QLIKE lower\\(DM sig.)} & \shortstack{Sharpe higher\\(HAC sig.)} \\
\midrule
per-bar linear & 3 & 0.0937 & 0.0942 & 0.0946 & 1.88 & 1.88 & 1.37 & 2 (0) & 2 (0) & 2 (0) & 3 (0) \\
pooled linear & 3 & 0.1038 & 0.1021 & 0.1065 & 1.23 & 1.44 & 0.89 & 1 (0) & 1 (0) & 3 (0) & 3 (0) \\
trees, shipped settings & 6 & 0.0924 & 0.0958 & 0.0986 & 1.47 & 1.50 & 0.93 & 6 (3) & 3 (1) & 6 (5) & 6 (0) \\
trees, random search & 12 & 0.0999 & 0.0995 & 0.1021 & 1.71 & 1.41 & 0.96 & 6 (3) & 9 (0) & 11 (0) & 11 (3) \\
trees, Optuna & 27 & 0.0974 & 0.1002 & 0.1010 & 1.30 & 1.23 & 0.71 & 22 (7) & 19 (1) & 26 (3) & 27 (6) \\
\midrule
\emph{all} & 51 & 0.0974 & 0.0994 & 0.1010 & 1.40 & 1.37 & 0.79 & 37 (13) & 34 (2) & 48 (8) & 50 (9) \\
\bottomrule
\end{tabular}
}
% \end{table}
% 
% \paragraph{A selection caveat.} The lean set was chosen after the importance study, whose
% refits cover the same period as these scores. Two checks bound what that costs. Ranking the
% series on the \flPreRefits{} refits trained before the trade period alone (their held-out
% tails, \flPreTailDays{} days from \flPreTailFirst{} to \flPreTailLast) picks the same five
% return and volume series: median ranks \flPreRankSumret{} for the signed return,
% \flPreRankSumabsret{} for the absolute return, \flPreRankSumvolume{} for E-mini volume,
% \flPreRankSumautocov{} for the autocovariance and \flPreRankSumretfour{} for the quartic return,
% against \flPreRankSumretthree{} for the cubed return, \flPreRankSumbipow{} for bipower
% variation and \flPreRankSumprettwo{} for the upside squared return. The cluster of the absolute,
% quartic and upside squared return and bipower variation already holds on the
% \flPreClusterRows{} rows from \flPreClusterFirst{} to \flPreClusterLast: its smallest
% within-cluster correlation is \flPreClusterMinCorrLagOne, \flPreClusterMinCorrLagFive,
% \flPreClusterMinCorrLagTwentyFive, \flPreClusterMinCorrLagOneTwentyFive,
% \flPreClusterMinCorrLagSixTwentyFive{} and \flPreClusterMinCorrLagThreeOneTwoFive{} at the
% lags from 1 to 3125 bars. The VIX level ranks \flRankMedVix{} on all tails and
% \flPreRankVix{} on the earlier ones: it is in the set by design, as the one
% implied-volatility series, and not by its rank. The set is therefore reproducible from
% importance computed only on refits trained before the trade period; what remains is that the
% selection rule itself was written after the full-sample results had been seen, and every
% comparison with \texttt{baseline} and \texttt{all\_features} carries that.
